\documentclass[10pt]{article}
\usepackage[a4paper,margin=19mm]{geometry}
\usepackage[T1]{fontenc}
\usepackage[utf8]{inputenc}
\usepackage{amsmath,amssymb,booktabs,longtable,array,xcolor}
\usepackage[hidelinks]{hyperref}
\title{Two-Sided Interval-Verified Optimality Bounds for the\\ PGLib-OPF v23.07 Typical AC Optimal Power Flow Cases}
\author{Puyi Sun\thanks{ORCID \href{https://orcid.org/0009-0004-1644-5300}{0009-0004-1644-5300}.}}
\date{\textbf{Computer-assisted proof candidate, not peer reviewed.}\\ Version 1.0, 28 October 2026}
\begin{document}\maketitle

\begin{abstract}
For each of the 26 typical-operating-condition (TYP) cases of the PGLib-OPF v23.07 benchmark library, with 3 to 2000
buses, we give a lower bound $L$ and an upper bound $U$ on the global optimum of the nonconvex AC optimal power flow
problem, together with machine-checkable certificates. The upper bounds come from Krawczyk interval-Newton enclosures of
exact power-flow solutions at which all operational limits are verified; the lower bounds come from sparse
semidefinite and sum-of-squares dual certificates whose positive semidefiniteness and polynomial residuals are checked in
interval arithmetic. All proof decisions use directed rounding on the unmodified case data. All 52 endpoints are accepted
by checkers written separately from the programs that produced the certificates. Thirteen of the 26 cases have a relative
gap of at most $10^{-6}$; the largest instances, case1951\_rte and case2000\_goc, have relative gaps of about
$8.1\times10^{-6}$ and $3.8\times10^{-4}$. The certificates, checkers and a one-command replay script accompany
this note. No claim of priority is made; the work was carried out with extensive use of AI systems, as detailed at the end.
\end{abstract}

\noindent\fcolorbox{black!40}{black!3}{\parbox{0.97\linewidth}{\textbf{Terminology.}
\emph{Rigorous / verified} means that every numerical proof obligation (positive semidefiniteness and residual bounds for $L$;
Krawczyk inclusion, every operational limit and the cost for $U$) is discharged by outward-rounded interval arithmetic on the
exact decimal data of the case file. In parts of the optimal power flow literature ``certified'' refers to a duality gap
below a solver tolerance; that weaker meaning is not used here. \emph{Independently checked} means accepted by checker
programs written separately from the certificate generators and importing nothing from them. It does \emph{not} mean
external human peer review or a formal proof in a proof assistant.}}

\section{Problem and claim}
Each case is the AC optimal power flow model of PGLib-OPF v23.07 \cite{pglib} with the data of its \texttt{.m} file taken as
exact decimals: complex bus voltages; nodal active and reactive power balance with shunts and $\pi$-model branches (tap
ratios and phase shifts); generator active and reactive power limits; voltage magnitude limits; apparent-power limits at
both ends of every rated branch; branch phase-angle-difference limits; and convex quadratic generation cost. Writing $K$
for the feasible set and $f$ for the cost, each case comes with a certificate pair proving
\begin{equation}
L \;\le\; f^\ast := \min_{x\in K} f(x) \;\le\; U .
\end{equation}
Both certificates of a case record the SHA-256 digest of the same case file. Other model variants (for example
current-magnitude thermal limits) and other PGLib versions are not covered.

\section{Upper bounds: Krawczyk enclosures}
An upper certificate stores an approximate operating point $\tilde z$ (voltages and free generator outputs) and a radius.
The checker rebuilds the power-flow equations $F(z)=0$ from the case file and forms the Krawczyk operator \cite{krawczyk}
\[
\mathcal K(\tilde z,Z)=\tilde z-C\,F(\tilde z)+\bigl(I-C\,F'(Z)\bigr)(Z-\tilde z),
\]
where the preconditioner $C$ is an approximate inverse proposed in floating point (sparse rows for large cases) and then
used as exact dyadic numbers. If $\mathcal K(\tilde z,Z)$ lies in the interior of the box $Z$, the box contains an exact
solution of $F=0$. The checker then evaluates every operational limit and the cost over the whole box in interval
arithmetic; if all limits hold, the enclosed exact solution is feasible and its cost is at most the upper end of the cost
interval. Setpoints fixed in the data (equal lower and upper limits) are compared as exact decimals. For the 2000-bus case
the system has 4000 real unknowns; one check took about 30--40 minutes and 160--210\,MB of memory on a two-core Linux
machine.

\section{Lower bounds: sparse dual certificates}
A lower certificate is a polynomial identity in the real voltage coordinates $x=(e,f)$,
\begin{equation}
f(x)-t=\sum_i \sigma_i(x)+\sum_j \tau_j(x)\,g_j(x)+\sum_k \eta_k(x)\,h_k(x)+R(x),
\end{equation}
where $g_j\ge 0$ and $h_k=0$ are the constraints, the $\sigma_i$ and the localizing multipliers $\tau_j$ are sums of squares
given by Gram matrices on the cliques of a chordal extension of the network graph (order one on most cliques and order two
on selected ones), the $\eta_k$ are polynomial multipliers, and $R$ is the residual. The checker proves each Gram matrix
positive semidefinite after a small diagonal shift by an interval $LDL^\top$ factorization (floating point only proposes the
shift), and subtracts from $t$ the worst-case effect of all shifts and of the residual over the voltage box. Bus costs
enter exactly, or through affine minorants $\lambda P+\psi(\lambda)$ whose constants are computed exactly. Nine cases use a chordal Shor-level (Lagrangian) certificate in a
dedicated format with its own checker; the remaining seventeen use the sparse sum-of-squares format.

\section{Results}
Table~\ref{tab:results} lists the bounds. Each value is the endpoint declared in the certificate, which the checker must
prove; $L$ is rounded down and $U$ up to six decimals, and the relative gap $(U-L)/U$ is computed from the full-precision
values. The widest gaps occur for case1803\_snem (3.6\%), case1888\_rte (1.9\%), case162\_ieee\_dtc
(1.4\%) and case240\_pserc (1.4\%). They describe the strength of the certificates available here, not a demonstrated limit
of any relaxation family; for case1803\_snem the lower certificate is chordal Shor-level, not an order-two certificate.

{\small\begin{longtable}{@{}l r r r r l@{}}\caption{Bounds for the 26 PGLib-OPF v23.07 TYP cases (\$/h).}\label{tab:results}\\\toprule
case & buses & $L$ & $U$ & $(U-L)/U$ & lower certificate\\\midrule\endfirsthead\toprule case & buses & $L$ & $U$ & $(U-L)/U$ & lower certificate\\\midrule\endhead
case3\_lmbd & 3 & 5812.642962 & 5812.643005 & 7.33e-09 & sparse SOS\\
case5\_pjm & 5 & 17551.890711 & 17551.891094 & 2.17e-08 & sparse SOS\\
case14\_ieee & 14 & 2178.080425 & 2178.080449 & 1.02e-08 & chordal Shor\\
case24\_ieee\_rts & 24 & 63352.202444 & 63352.203674 & 1.94e-08 & chordal Shor\\
case30\_as & 30 & 803.127218 & 803.127313 & 1.17e-07 & chordal Shor\\
case30\_ieee & 30 & 8208.515455 & 8208.515549 & 1.13e-08 & chordal Shor\\
case39\_epri & 39 & 138415.554411 & 138415.563671 & 6.69e-08 & sparse SOS\\
case57\_ieee & 57 & 37589.336873 & 37589.338384 & 4.02e-08 & sparse SOS\\
case60\_c & 60 & 92693.623915 & 92693.670117 & 4.98e-07 & sparse SOS\\
case73\_ieee\_rts & 73 & 189764.081348 & 189764.084939 & 1.89e-08 & chordal Shor\\
case89\_pegase & 89 & 106968.658190 & 107285.675532 & 2.95e-03 & chordal Shor\\
case118\_ieee & 118 & 97213.585027 & 97213.607809 & 2.34e-07 & sparse SOS\\
case162\_ieee\_dtc & 162 & 106512.858738 & 108075.648889 & 1.45e-02 & sparse SOS\\
case179\_goc & 179 & 754128.038771 & 754266.421037 & 1.83e-04 & sparse SOS\\
case197\_snem & 197 & 1.501257 & 1.501747 & 3.25e-04 & chordal Shor\\
case200\_activ & 200 & 27557.570830 & 27557.571313 & 1.75e-08 & chordal Shor\\
case240\_pserc & 240 & 3283231.188888 & 3329670.118665 & 1.39e-02 & sparse SOS\\
case300\_ieee & 300 & 565219.433978 & 565220.008801 & 1.02e-06 & sparse SOS\\
case500\_goc & 500 & 454945.910696 & 454945.987308 & 1.68e-07 & sparse SOS\\
case588\_sdet & 588 & 311343.651681 & 313139.784598 & 5.74e-03 & sparse SOS\\
case793\_goc & 793 & 258581.712068 & 260197.852143 & 6.21e-03 & sparse SOS\\
case1354\_pegase & 1354 & 1257170.185668 & 1258844.004024 & 1.33e-03 & sparse SOS\\
case1803\_snem & 1803 & 94825.222708 & 98334.547531 & 3.57e-02 & chordal Shor\\
case1888\_rte & 1888 & 1375653.865114 & 1402531.010635 & 1.92e-02 & sparse SOS\\
case1951\_rte & 1951 & 2085565.057043 & 2085581.868859 & 8.06e-06 & sparse SOS\\
case2000\_goc & 2000 & 973063.826000 & 973432.478973 & 3.79e-04 & sparse SOS\\
\bottomrule\end{longtable}
}

\section{Independent checking and replay}
Three checker programs, one for upper certificates and one for each lower-certificate format, were written separately from
the generators; each reads only the raw case file and one certificate, uses the \texttt{mpmath} interval library for every
decision, and uses floating point only to propose numbers that are then verified. All 52 endpoints were accepted by these
checkers on a Linux machine, and again, all 52, on the author's own computer (WSL2 Ubuntu, Python 3.14) with byte-identical inputs;
the smaller cases were also replayed in a separate sandbox, and many endpoints on a Windows machine. Deliberately corrupted
certificates (zeroed or negated Gram blocks, negated multipliers, altered cost pieces, declared bounds raised above the
proved value, wrong case-file digest, removed equality multipliers, Krawczyk centres moved, setpoints pushed out of
limits) were all rejected. The accompanying \texttt{replay.py} verifies the SHA-256 manifest of the release, checks that the
declared values in its index match the certificates, and re-runs every checker.

\section{Related work and position}
Oustry, D'Ambrosio, Liberti and Ruiz \cite{oustry} computed certified SDP lower bounds (via a Gershgorin-type eigenvalue
bound on rational dual data) for AC-OPF instances with up to 6515 buses (PGLib v21.07, with current-magnitude limits);
their upper bounds are local-solver (IPOPT) values without a feasibility proof. Gopinath et al.\ \cite{gopinath} closed or nearly closed many PGLib gaps with SDP-based bound
tightening and valid inequalities, using floating-point solvers. Sparse moment-sum-of-squares hierarchies certify global
optimality numerically on large AC-OPF instances \cite{cstssos,wangopf,lefranc}. Recent work strengthens relaxations with
offline bound tightening \cite{an}, lifted nonlinear cuts \cite{bugosen} and enhanced SDP models \cite{enhancedsdp}; robust
AC-OPF has been solved to small solver-tolerance gaps by spatial branch and bound \cite{yin}. Interval methods for power
flow feasibility are classical. We make no priority claim. What this note adds is a fixed, machine-replayable artifact in
which both ends of the optimality interval are verified with directed rounding, for all 26 TYP cases up to 2000 buses, and
re-checked by separately written checkers. We could not read the full text of \cite{an}; its abstract reports
numerical relaxation gaps only.

\section{Limitations}
\begin{enumerate}
\item The results rest on the correctness of the checker programs and of the interval library; separate authorship of
checkers and generators reduces, but does not remove, the risk of a shared error. There has been no external review.
\item The bounds concern the PGLib v23.07 TYP model exactly as written; they say nothing about other variants, other data or
real-world operation.
\item Several gaps are wide (Section~4); narrowing them needs stronger relaxations than those used here.
\end{enumerate}

\section*{AI assistance and responsibility}
This work was carried out with AI systems under the direction of the author, who chose the problem, set the acceptance
rules, decided which results to keep, operated the machines and is responsible for this release. The certificate-generating
programs were written mainly by Claude (Anthropic), with one SDP-assembly module by Codex (OpenAI); certificates were
generated by Claude and, for some lower certificates, by Codex. The independent checkers were written by GPT (OpenAI) and
their source was reviewed by Claude; they were run by Claude, by GPT, by Codex and on the author's computer as stated in
Section~5. Mathematical review of the certificate formats was done by GPT and Claude, reviewing each other. This text was
drafted by GPT and revised by Claude. \emph{The author did not personally verify the mathematics line by line}; the claims
rest on the certificates and checkers, which anyone can re-run. AI systems are not listed as authors.

\section*{Data and code}
Certificates, checkers, case files and the replay script are in the accompanying repository
(https://github.com/PYCSON/acopf-bounds-proofs). Code: MIT licence. Certificates and text: CC BY 4.0. Case files: PGLib-OPF v23.07 \cite{pglib},
CC BY 4.0, unmodified.

\bibliographystyle{plain}
\bibliography{references}
\end{document}
